\section*{ЗАКЛЮЧЕНИЕ}
\label{sec:conclusion}
\addcontentsline{toc}{section}{ЗАКЛЮЧЕНИЕ}

Настоящата дипломна работа представи проектирането, реализацията и анализа на библиотека за обектно-релационно картографиране по време на компилация за C++. Логическият модел на данните и структурата на заявките се описват чрез статично типизирани конструкции; релационните и нерелационните системи за съхранение се интегрират чрез формален договор на конектора и маркери за способности (capability tags); а асинхронното изпълнение е изградено върху корутини. Архитектурата е валидирана чрез 271 единични и интеграционни теста и емпирична оценка на асинхронния модел.

Основните приноси са четири. Първо, демонстрирано е, че практически полезна compile-time ORM библиотека може да бъде реализирана в стандартен C++ без външна кодогенерация. Второ, изяснена е връзката между C++ модела и физическата организация на данните: моделът в C++ служи като първичен носител на логическата схема, а конкретната система за съхранение определя материализацията. Трето, формализиран е договорът на конектора, който позволява надграждане с нови backend-и без разрушаване на гаранциите на интерфейса. Четвърто, изграден е асинхронен слой върху корутините на C++20 с честно разграничение между универсалното изнасяне в обединение от нишки и интеграцията през неблокиращите интерфейси на драйверите --- емпиричната оценка показва ускорение до 23$\times$ в режим, ограничен от вход-изход.

Работата има ясни граници: общото междинно представяне на заявката не може да бъде интерпретирано еднакво върху всеки модел за съхранение, а зрелостта на интеграцията през неблокиращите интерфейси варира според това какво позволява съответният драйвер. Бъдещото развитие следва пряко от вече реализираните решения --- по-зряло live покритие, автоматизация на схемата, разширяване на неблокиращата интеграция, интегриране на отражение (reflection) на C++26 и по-фина таксономия на способностите.

Комбинацията от статичен обектен модел, типизирано междинно представяне, договор на конектора, основан на признаци (traits), асинхронен слой над корутини и систематична верификация показва устойчив път към библиотеки за достъп до данни, които са едновременно изразителни, безопасни и разширяеми. Този подход не е теоретично упражнение, а реална инженерна алтернатива за системите, изискващи строгост, яснота и разширяемост при работа с множество модели за съхранение.

В по-широк смисъл, работата показва, че \emph{типовата система на C++ е достатъчно изразителна, за да формализира не само описанието на данни и алгоритми, но и архитектурните граници, допустимостта на операциите и прехода между логически и физически слоеве}. Обектно-релационното картографиране, дълго време разглеждано като неестествена употреба на C++, се оказва не само осъществимо без външна кодогенерация, а и научно интересен случай: то изисква едновременно метапрограмиране на ниво типове, корутинна теория за асинхронното изпълнение и архитектурна дисциплина за разделяне на логическата форма от физическата материализация. Тази тройна връзка между типове, изпълнителни модели и архитектурни граници очертава жанр от библиотеки, в които компилаторът поема не само ролята на оптимизатор, а и на първи проверяващ на коректността --- и в който C++ показва своя зрелостен момент като език за изграждане на инфраструктурни системи с висока степен на формалност.
